\subsection{The singular-parameter locus of the characteristic-two Moore matrix}

This subsection concerns an explicit polynomial matrix.  It makes no
identification with a Cartier tensor.  Let \(k\) be algebraically closed
of characteristic two, fix a primitive cube root \(\zeta\), and put
\(G=\mathbf F_3^2\).  For \(c=(c_v)_{v\in G}\) and independent
variables \(x=(x_v)_{v\in G}\), define
\[
 M_c(x)_{a,b}=c_{a-b}x_{a+b},\qquad
 D_c(x)=\det M_c(x),\qquad a,b\in G.                 \tag{DE-MC1}
\]
For \(h\in G\) and \(\ell\in G^\vee\), let
\[
 V_{h,\ell}=\{c:\ c_h=0,\quad
 c_{h+v}=\zeta^{\ell(v)}c_{h-v}\text{ for every }v\in G\}.
                                                               \tag{DE-MC2}
\]
Each is a four-dimensional vector space.  Relations for one member
of each opposite nonzero pair give its four independent equations
in addition to \(c_h=0\).

\begin{lemma}[Exact support and coefficient reductions]
\label{lem:DE-Moore-support}
\status{Candidate; complete finite symbolic certificate and proof supplied.}
Suppose \(c\ne0\) and \(D_c\) is identically zero.  Its nonzero
coordinate support is symmetric about a missing coordinate.  The
possible supports form four orbits under affine changes of the
\(G\)-indices:
\[
\begin{array}{c|c|c}
\text{size}&\text{representative support}&\text{number of supports}\\\hline
2&\{(0,1),(0,2)\}&36\\
4&\{(0,1),(0,2),(1,0),(2,0)\}&54\\
6&\{(i,j):i\ne0\}&12\\
8&G\setminus\{0\}&9.
\end{array}                                                     \tag{DE-MC3}
\]
The following are actual coefficients of the restricted determinant,
up to their indicated nonzero monomial factors.
On the two-coordinate support the entire determinant is
\[
 (c_{01}^3+c_{02}^3)^3\prod_{v\in G}x_v.                         \tag{DE-MC4}
\]
On the four-coordinate support, two coefficients are
\[
 (c_{01}^3+c_{02}^3)c_{10}^3c_{20}^3,
 \qquad (c_{10}^3+c_{20}^3)c_{01}^3c_{02}^3.                     \tag{DE-MC5}
\]
On the six-coordinate support, put \(A_i=c_{(1,i)}\) and
\(B_i=c_{(2,i)}\), for \(i\in\mathbf F_3\).  Five coefficients are
\[
 \left(\prod_i A_i+\prod_i B_i\right)^3,
 \qquad\left(\sum_i A_i^3+\sum_i B_i^3\right)^3,
 \qquad\prod_i(A_i^3+B_{i+r}^3)\quad(r\in\mathbf F_3).
                                                               \tag{DE-MC6}
\]
On the eight-coordinate support, four coefficient binomials force
\(c_v^3=c_{-v}^3\).  After writing
\[
 c_v=\zeta^{r_v}d_v,\quad c_{-v}=d_v,
 \qquad v\in\{(0,1),(1,0),(1,1),(1,2)\},\quad d_v\ne0,
                                                               \tag{DE-MC7}
\]
the only phase quadruples compatible with an identically zero
determinant are \((r_v)_v=(\ell(v))_v\) for \(\ell\in G^\vee\).
\end{lemma}
\begin{proof}
Here all support and phase reductions are finite symbolic identities,
for which we give the exact reconstruction procedure and retain
every exclusion witness.  Index \(G\) lexicographically.  For each
of the \(9!\) permutations \(\sigma\) of its indices, form
\[
 u_g(\sigma)=\#\{a:a-\sigma(a)=g\},\qquad
 v_g(\sigma)=\#\{a:a+\sigma(a)=g\}.
\]
The determinant in characteristic two is the sum of
\(c^{u(\sigma)}x^{v(\sigma)}\); identical monomials cancel in pairs.
This produces exactly \(55{,}296\) monomials grouped into
\(2{,}080\) nonzero \(x\)-coefficient polynomials.
For each of the \(511\) nonempty subsets \(S\subset G\), discard
all monomials using a \(c\)-coordinate outside \(S\).  In exactly
\(400\) supports, one resulting coefficient consists of a single
monomial.  Its value cannot vanish on the torus with support exactly
\(S\), so that support is impossible.  The retained certificate
lists both exponent vectors for each such witness.  The remaining
\(111\) subsets are exactly those in (DE-MC3); for every one the
certificate gives an affine map from the displayed representative.
Thus no test at finitely many values of the nonzero coordinates
enters this support exclusion.

For completeness, the four orbit counts also have a direct
description.  A two-point support has its unique midpoint.  A
four-point support is the union of two opposite pairs about its
missing center; that center is unique, since two different centers
would give a translation of order three preserving four points.
There are \(9\binom42=54\) such supports.  A six-point support is
the complement of an affine line, giving twelve choices.  An
eight-point support is the complement of one point.  The affine
group is transitive on each of these four types.  Its compatibility
with determinant vanishing is the literal reindexing identity
\[
 M_c(x)_{Ta+h,Tb}=c_{T(a-b)+h}\,x_{T(a+b)+h}
 \qquad(T\in\operatorname{GL}_2(\mathbf F_3)).                   \tag{DE-MC8}
\]
Rows, columns and the independent \(x\)-variables are only
permuted.

The same reconstructed coefficient table gives (DE-MC4)--(DE-MC6)
by multiplication of their displayed polynomial factors in
\(\mathbf F_2[c]\).  On the eight-coordinate chart the coefficient
of \(x_0^7x_vx_{-v}\), after setting \(c_0=0\), is
\[
 (c_v^3+c_{-v}^3)\prod_{u\notin\{0,v,-v\}}c_u.                 \tag{DE-MC9}
\]
Its remaining six factors are nonzero.  Therefore each opposite
ratio is a cube root of unity, giving exactly the \(3^4=81\)
phase choices in (DE-MC7).  Substitute every such choice into
every coefficient while retaining the four independent variables
\(d_v\).  Exactly the nine character choices make all coefficients
zero as polynomials in these variables.  For each of the other
seventy-two choices, the certificate retains a coefficient equal
to a single monomial in the \(d_v\), with nonzero coefficient in
\(\mathbf F_4\).  It is nonzero everywhere on this chart.
This exhaustive root-of-unity substitution proves the asserted
necessity over \(k\); it is not an extrapolation from its
\(\mathbf F_4\)-valued points.  The complete standalone program
and exact witnesses are specified in the final remark.
\end{proof}

\begin{lemma}[Three cyclic matching products]
\label{lem:DE-Moore-six-matching}
\status{Candidate; complete algebraic proof supplied.}
Let \(a_i,b_i\), \(i\in\mathbf F_3\), be nonzero elements of a
field of characteristic two.  Suppose
\[
 \sum_i a_i=\sum_i b_i,\qquad
 \prod_i a_i=\prod_i b_i,\qquad
 \prod_i(a_i+b_{i+r})=0\quad(r\in\mathbf F_3).
                                                               \tag{DE-MC10}
\]
Then there is \(r_0\in\mathbf F_3\) such that
\(a_i=b_{r_0-i}\) for all \(i\).
\end{lemma}
\begin{proof}
For each \(r\), choose \(i_r\) with \(a_{i_r}=b_{i_r+r}\).
View these as three edges between the three \(a\)-indices and
the three \(b\)-indices, one edge for each cyclic difference.
If the row indices \(i_r\) are distinct, the three column
indices have sum zero in \(\mathbf F_3\).  They are therefore
either all distinct or all equal.  When both sets of indices
are distinct, the matching is a permutation of the three-element
set.  Every such permutation is \(i\mapsto i+t\) or
\(i\mapsto-i+t\).  The first has a constant difference,
contrary to the three selected differences.  The second is
the required reflection matching.

If exactly two row indices occur, exactly two column indices
occur: writing out the unused difference verifies this immediately.
The three distinct edges in this two-by-two rectangle connect
all four vertices.  Their four values are consequently equal.
The equality of the two sums then identifies the unused row value
with the unused column value.  The two bijections between the
equal row pair and the equal column pair give one translation
permutation and one reflection permutation, so a reflection
matching is again available.

The remaining cases are a three-edge star on either side.
Three values on one side and one value on the other side are
then a common nonzero value \(u\).  Equality of sums makes
the other two values equal, say to \(v\).  Equality of products
gives \(uv^2=u^3\), so \(v^2=u^2\) and hence \(v=u\).
All six values are equal.  Any reflection works.
These cases exhaust the three edge choices.  The retained
certificate independently records all twenty-seven choices,
including the eighteen two-by-two cases and the six stars.
\end{proof}

\begin{theorem}[All identically singular Moore parameters]
\label{thm:DE-Moore-singular-parameters}
\status{Candidate; complete proof and universal coefficient certificate supplied.}
Set-theoretically in \(\mathbf P^8_k\),
\[
 \{[c]:D_c(x)\equiv0\}
       =\bigcup_{h\in G,\ \ell\in G^\vee}\mathbf P(V_{h,\ell}).
                                                               \tag{DE-MC11}
\]
The right side consists of eighty-one distinct projective
three-spaces.
\end{theorem}
\begin{proof}
First suppose \(c\in V_{h,\ell}\).  Permute the rows of
\(M_c\) by \(a\mapsto a+h\), and multiply row \(a\) by
\(\zeta^{-\ell(a)}\).  The resulting entry is
\[
 \zeta^{-\ell(a)}c_{h+a-b}x_{h+a+b}.
\]
The defining relation (DE-MC2) makes this matrix symmetric,
and its diagonal is zero because \(c_h=0\).  It is therefore
alternating of odd size and has zero determinant.  This argument
is a polynomial identity, also when some of its four free
coordinates vanish.

Conversely, assume \(c\ne0\) and \(D_c\equiv0\).
The support lemma and affine covariance reduce to one of the
four displayed supports.  On a two-point support, (DE-MC4)
gives \(c_{01}^3=c_{02}^3\).  Its opposite-coordinate ratio
is a cube root of unity, which is the restriction of a character
of \(G\).  Thus \(c\) lies in a space (DE-MC2).
On the four-point support, (DE-MC5) gives the two corresponding
cube-ratio identities.  The two directions are independent,
so their chosen cube-root ratios extend to a unique character
of \(G\), with the same conclusion.  On the eight-point support,
the final assertion of the support lemma gives the conclusion
directly.

Consider the six-point support.  From (DE-MC6) and the absence
of nonzero nilpotents in a field, we have
\[
 \prod_iA_i=\prod_iB_i,\qquad
 \sum_iA_i^3=\sum_iB_i^3,\qquad
 \prod_i(A_i^3+B_{i+r}^3)=0\quad(r\in\mathbf F_3).
                                                               \tag{DE-MC12}
\]
Apply the matching lemma to \(a_i=A_i^3\), \(b_i=B_i^3\);
their product equality follows by cubing the first equality.
There is \(r_0\) with \(A_i^3=B_{r_0-i}^3\).
Write
\[
 \eta_i=A_i/B_{r_0-i}=\zeta^{t_i},\qquad t_i\in\mathbf F_3.
\]
The first equality in (DE-MC12) gives
\(\prod_i\eta_i=1\), or \(t_0+t_1+t_2=0\).
Every such triple is of the form
\(t_i=\alpha+\beta i\): take
\(\alpha=t_0\), \(\beta=t_1-t_0\), and the third equality
follows from its zero sum.

Choose \(h=(0,h_2)\) with \(2h_2=r_0\), and choose the
linear character \(\ell\) whose values satisfy
\[
 \ell(1,i-h_2)=\alpha+\beta i.
\]
For instance its two coefficients are
\((\alpha+\beta h_2,\beta)\).  The equations
\(A_i=\zeta^{\ell(1,i-h_2)}B_{2h_2-i}\) are exactly the
opposite-pair equations about \(h\).  The three coordinates
on the missing affine line are all zero, so their remaining
opposite-pair equation and \(c_h=0\) hold as well.
Thus \(c\in V_{h,\ell}\), completing the necessity.

Finally, each \(V_{h,\ell}\) has four independent free
coordinates, one per opposite pair.  Its only identically zero
individual coordinate is \(c_h\), so spaces with different
centers are distinct.  For a fixed center, two different linear
characters differ on some direction, hence give different
pair equations.  This proves the asserted count and dimension.
\end{proof}

\begin{remark}[Exact certificate and the geometric limitation]
\label{rem:DE-Moore-classification-scope}
\status{Candidate scope; no Cartier or fixed-surface rank implication asserted.}
The standalone program
\path{code/obstructions_crystals/verify_Moore_singular_parameter_classification.py}
reconstructs all \(362{,}880\) determinant terms before cancellation.
It checks all \(511\) supports and retains all \(400\) single-monomial
exclusions; checks \(432\) affine index maps and supplies maps for
all \(111\) surviving supports; reconstructs the eight factored
coefficients (DE-MC4)--(DE-MC6); checks all \(27\) matching-edge
choices and the four exact coefficient locators (DE-MC9); substitutes all \(81\) eight-support phase patterns into
every universal coefficient, retaining the \(72\) noncharacter
monomial witnesses; and checks \(6{,}561\) symbolic entries in the
alternating row-scalings of all \(81\) spaces.  Its completed author
run passed in \(3.2868060610053362\) seconds at
01:26:52 UTC on 20 November 2026.  No execution failed.
The entire code, output, standard streams, earlier partial probes
and all rejected shortcuts are retained.  These are exhaustive
finite symbolic reductions of the stated universal matrix,
not finite-field estimates of an unspecified geometric locus.

The theorem concerns the set of geometric parameters; it does not
identify the scheme structure of the coefficient ideal.  In
particular, cubes in some of the displayed coefficients are not
silently replaced by equations at the scheme-theoretic level.
No theta embedding is identified with the parameter vector \(c\),
and no actual twisted Cartier cross-pairing is identified with
(DE-MC1).  Those typed geometric identifications remain necessary
before this result can imply a rank statement for a fixed abelian
surface.  No main rank or monodromy case is asserted here.

The earlier forty-basis coordinate-product shortcut is false:
the retained exact parameter \((1,1,1,1,0,0,0,0,0)\) satisfies
all forty products but has a nonzero full determinant.  The
present proof uses the full mixed coefficient polynomials and
does not reuse that rejected inference.  No external
classification theorem is imported, and no historical-priority
claim is made for the explicit algebraic classification.
\end{remark}
